% !TEX root = main.tex
\chapter{Introduction} 
\label{chap_introduction}
\section{Introduction \& Motivation} 
\label{sec_introduction_motivation}
\begin{center}
   \textit{Aim of this report is to identify the possible procurement paths of Tritium, one of the two educts necessary for a nuclear fusion reaction. For each path the chemical background will be identified and each pathway shall be characterized.}\\
\end{center}

Fusion technology is on the upswing in recent years: Research and Development projects, privately and
publicly funded, have sprouted all over the globe. There are 103 fusion devices in operation at the 
moment, with an additional 19 under construction and 60 planned projects.
\footnote{Website last accessed on 07-05-2026} \cite{IAEA_FFDB_2025}
However, it is central for the understanding of the technology and its global standing to emphasise that 
there is no commercial-ready fusion plant implemented today. Its commercial feasibility still needs to 
be proven. \cite{wimmersNuclearPowerTechnology2026} %, pp. 99 \& 130

The nuclear fusion supply chain is immature and therefore sensitive. \cite{wimmersNuclearPowerTechnology2026} %pp. 115 \& 129
One of the reasons for this critical state of the supply chain is the small number of producers and vendors along the chain. For the supply chain to be further developed the commercial application might have to be proven first, thereby promoting the sector for investors and subsequently stabilizing the chain.  \cite{wimmersNuclearPowerTechnology2026} %pp. 115

A link in this supply chain, which is especially volatile is nuclear fusions fuel. More specifically the material needed to breed tritium, highly enriched lithium-6 isotopes.

The current state of affairs leads to either an global bottleneck of lithium-6 or total European dependency on imports of lithium-6 as a end of the chain product.

Global demand for some raw materials will increase distinctly in the near future. One of these materials is lithium, with estimates that its demand will rise up to 13 times until 2040. \cite{colucciCombinedAssessmentMaterial2025}

This thesis is organized as follows. First section \ref{sec_technological_background} gives the necessary background information and further talks about existing literature. At the end of this section the novelty of this work against the beforehand backdrop will be highlighted. 
\ref{xxx} lays out the methodology applied including the utilized approaches and mathematical formulas. 
\ref{} talks about scenarios. 
\ref{xxx} showcases the results of this work, with follow-up section \ref{xxx} summing op the whole thesis with some conclusory paragraphs.  